
This improvement is consistent with reduced contamination of the reference amplitude distribution by heterogeneous epoch composition. When blink-heavy and blink-free epochs were pooled across the concatenated recording, the large proportion of blink-free signal pulled the estimated threshold toward the baseline amplitude, causing the detector to fire too readily. The resulting error structure provides direct evidence for this mechanism: BLINKER-concat retained the highest pooled recall of 0.9627 but had a precision of 0.6182, with 627.9 false positives and 23.5 false negatives per session and an FP:FN ratio of 26.675. Its errors were therefore overwhelmingly false-positive-heavy rather than arising from missed blinks. This sensitivity is expected because ocular-artifact amplitude varies across subjects, electrodes, and recording sessions~\citep{croft2000removal,guttmann2019new,klein2013reliable,kleifges2017blinker}.
A single threshold derived from the full recording consequently depends on the relative contribution of blink-rich and blink-poor periods, whereas epoch screening restricts threshold estimation to signal segments that better represent the local background. The observed precision gain therefore reflects improved calibration to session-specific amplitude structure rather than reduced sensitivity to true blinks.
